% SPDX-License-Identifier: CC-BY-SA-4.0
% Author: Matthieu Perrin

\begingroup

\begin{exercice}[Exécutions d'un programme réparti]
  
  On considère trois processus $p_1,p_2,p_3$ dans un modèle asynchrone sans crash.
    
  \begin{minipage}[t]{.24\textwidth}
    \begin{algorithm}[H]
      \Process{$p_1$}{
        \nl \Send $m_1$ \To $p_3$\\
        \nl \Send $m_2$ \To $p_2$
      }
    \end{algorithm}
  \end{minipage}%
  %
  \begin{minipage}[t]{.37\textwidth}
    \begin{algorithm}[H]
      \setcounter{AlgoLine}{2}
      \Process{$p_2$}{
        \nl \Send $m_3$ \To $p_3$\;
      }
      \Upon{\Receive $m_2$ \From $p_1$ \At $p_2$}{
        \nl \Send $m_4$ \To $p_3$
      }
    \end{algorithm}
  \end{minipage}%
  %
  \begin{minipage}[t]{.37\textwidth}
    \begin{algorithm}[H]
      \setcounter{AlgoLine}{4}
      \Upon{\Receive $m_k$ \From $p_j$ \At $p_3$}{
        \nl $\Print(k)$
      }
    \end{algorithm}
  \end{minipage}

  \begin{questions}
  \item Décrire la machine à états du programme ci-dessus. 
  \end{questions}

  \begin{questions}
  \item Décrire toutes les exécutions causales possibles. 
  \end{questions}

\end{exercice}

\endgroup
\endinput
